% Chapter 4: Submersions, Immersions, and Embeddings
% Lee, \emph{Introduction to Smooth Manifolds} (2nd ed., GTM 218).
% Generated from the section extraction; nodes carry only Lean that is
% verified sorry-free. See ../../UPSTREAM_LEAN_AUDIT.md.

\section{Local Diffeomorphisms}

\begin{definition}
\label{def:4-22-extra-1}
\lean{Manifold.isSmoothLocalDiffeomorphism_iff_forall_isLocalDiffeomorphAt}
\leanok


For smooth manifolds $M$ and $N$ with or without boundary, a map $F:M\to N$ is a local diffeomorphism if every $p\in M$ has a neighborhood $U$ such that $F(U)$ is open in $N$ and $F|_U:U\to F(U)$ is a diffeomorphism.
\end{definition}

\begin{theorem}[Inverse Function Theorem for Manifolds]
\label{thm:4-5}
\lean{exists_connected_open_neighborhoods_diffeomorph_of_mfderiv_isInvertible}
\leanok


\dcref{ch4:4.5}
Suppose $M$ and $N$ are smooth manifolds and $F:M\to N$ is smooth. If $p\in M$ and $dF_p$ is invertible, then there are connected neighborhoods $U_0$ of $p$ and $V_0$ of $F(p)$ such that $F|_{U_0}:U_0\to V_0$ is a diffeomorphism.
\end{theorem}

\begin{proof}
\leanok
\sketch
Bijectivity of $dF_p$ gives $\dim M=\dim N=n$. Choose charts $(U,\varphi)$ and $(V,\psi)$ centered at $p$ and $F(p)$ with $F(U)\subseteq V$. The coordinate map $\widehat F=\psi\circ F\circ\varphi^{-1}$ satisfies $\widehat F(0)=0$ and has nonsingular differential at $0$, since
\[
d\widehat F_0=d\psi_{F(p)}\circ dF_p\circ d(\varphi^{-1})_0.
\]
The Euclidean inverse function theorem yields connected open sets $\widehat U_0\subseteq\varphi(U)$ and $\widehat V_0\subseteq\psi(V)$ containing $0$ such that $\widehat F|_{\widehat U_0}$ is a diffeomorphism onto $\widehat V_0$. Their inverse images $U_0=\varphi^{-1}(\widehat U_0)$ and $V_0=\psi^{-1}(\widehat V_0)$ have the required properties.
\end{proof}

\begin{proposition}[Detecting Smoothness by Local Diffeomorphisms]
\label{exer:4-10}
\lean{smooth_iff_comp_right_of_surjective_isLocalDiffeomorph}
\leanok


\dcref{ch4:4.10}
Let $M,N,P$ be smooth manifolds with or without boundary and let $F:M\to N$ be a local diffeomorphism. If $G:P\to M$ is continuous, then $G$ is smooth if and only if $F\circ G$ is smooth. If $F$ is also surjective and $G:N\to P$ is any map, then $G$ is smooth if and only if $G\circ F$ is smooth.
\end{proposition}

\section{Embeddings}

\begin{definition}
\label{def:4-24-extra-1}
\lean{Manifold.IsSmoothEmbedding}

\mathlibok
A smooth embedding $F:M\to N$ between smooth manifolds with or without boundary is a smooth immersion that is also a topological embedding, namely a homeomorphism from $M$ onto $F(M)$ with the subspace topology.
\end{definition}

\begin{example}[The Figure-Eight Curve]
\label{ex:4-19}
\lean{figureEightCurve_closed_range_eq_restricted_range}
\leanok


\dcref{ch4:4.19}
The map $\beta:(-\pi,\pi)\to\mathbb R^2$ given by
\[
\beta(t)=(\sin 2t,\sin t)
\]
is an injective smooth immersion but not a topological embedding.

\sketch Its derivative $(2\cos 2t,\cos t)$ never vanishes, while its image is compact and its domain is not.
\end{example}

\begin{lemma}[Dirichlet's Approximation Theorem]
\label{lem:4-21}
\lean{dirichlet_approximation_theorem}
\leanok


\dcref{ch4:4.21}
For every $\alpha\in\mathbb R$ and positive integer $N$, there are integers $n,m$ with $1\le n\le N$ such that $|n\alpha-m|<1/N$.
\end{lemma}

\begin{proof}
\leanok
\sketch
Set $f(x)=x-\lfloor x\rfloor$. The $N+1$ values $f(i\alpha)$ for $0\le i\le N$ lie in $[0,1)$, so two, say with $0\le i<j\le N$, lie in the same one of the $N$ intervals $[k/N,(k+1)/N)$. Hence $|f(j\alpha)-f(i\alpha)|<1/N$. Taking $n=j-i$ and $m=\lfloor j\alpha\rfloor-\lfloor i\alpha\rfloor$ proves the claim.
\end{proof}


\section{Submersions}

\begin{definition}
\label{def:4-25-extra-1}
\lean{Function.RightInverse}

\mathlibok
For a continuous map $\pi:M\to N$, a section is a continuous map $\sigma:N\to M$ satisfying $\pi\circ\sigma=\operatorname{Id}_N$. A local section is a continuous map $\sigma:U\to M$ on an open set $U\subseteq N$ satisfying $\pi\circ\sigma=\operatorname{Id}_U$.
\end{definition}

\section{Smooth Covering Maps}

\begin{definition}
\label{def:4-26-extra-1}
\lean{Manifold.refl_isLocalDiffeomorph_id}
\leanok


A topological covering map is a surjective continuous map $\pi:E\to M$ between connected, locally path-connected spaces such that every $q\in M$ has an evenly covered neighborhood $U$: each component of $\pi^{-1}(U)$ maps homeomorphically onto $U$. For connected smooth manifolds $E$ and $M$ with or without boundary, a smooth covering map is a smooth surjection such that every $q\in M$ has a neighborhood $U$ for which each component of $\pi^{-1}(U)$ maps diffeomorphically onto $U$. The manifold $M$ is the base, $E$ a covering manifold, and a simply connected $E$ is a universal covering manifold. Thus a smooth covering map is a topological covering map with the additional diffeomorphism requirement.
\end{definition}

\begin{proposition}[Topologically Evenly Covered Neighborhoods]
\label{exer:4-39}
\lean{topological_evenly_covered_component_isLocalDiffeomorphOn_and_bijOn}
\leanok


\dcref{ch4:4.39}
If $\pi:E\to M$ is a smooth covering map and $U\subseteq M$ is evenly covered
in the topological sense, each component of $\pi^{-1}(U)$ maps
diffeomorphically onto $U$. Thus the topological and smooth notions of evenly
covered agree.
\end{proposition}

\begin{proposition}
\label{prop:4-46}
\lean{isSmoothCoveringMap_of_proper_localDiffeomorph}
\leanok


\dcref{ch4:4.46}
Suppose $E$ and $M$ are nonempty connected smooth manifolds with or without boundary. A proper local diffeomorphism $\pi:E\to M$ is a smooth covering map.
\end{proposition}

\begin{proof}
\leanok
\sketch
The local diffeomorphism $\pi$ is open and properness makes it closed, so its nonempty image is both open and closed in connected $M$; hence $\pi$ is surjective. For $q\in M$, the fiber $\pi^{-1}(q)$ is discrete by local injectivity and compact by properness, hence finite, say $\{p_1,\ldots,p_k\}$. Choose pairwise disjoint neighborhoods $V_i$ on which $\pi$ is a diffeomorphism onto open $U_i\ni q$, and set $U=\bigcap_iU_i$. With
\[
K=E\setminus(V_1\cup\cdots\cup V_k),
\]
Since $q\notin\pi(K)$, replace $U$ by $U\setminus\pi(K)$ and then by its connected component containing $q$. Then
\begin{equation}
U\subseteq U_i\quad\text{for every }i, \tag{4.8}
\end{equation}
and
\begin{equation}
\pi^{-1}(U)\subseteq V_1\cup\cdots\cup V_k. \tag{4.9}
\end{equation}
The sets $\widetilde V_i=\pi^{-1}(U)\cap V_i$ are disjoint, connected, open, cover $\pi^{-1}(U)$, and each maps diffeomorphically onto $U$. They are therefore exactly its components, so $U$ is evenly covered.
\end{proof}

\begin{exercise}
\label{exer:4-24}
\lean{origin_mem_closure_range_oscillatingCurveMap}
\leanok
\dcref{ch4:4.24}
\begin{itemize}
\item Exercise 4.24. Give an example of a smooth embedding that is neither an open map nor a closed map.
\end{itemize}
\end{exercise}

\begin{exercise}
\label{exer:4-4}
\lean{Manifold.HasConstantRank}
\leanok
\dcref{ch4:4.4}
Show that a composition of smooth submersions is a smooth submersion, and a composition of smooth immersions is a smooth immersion. Give a counterexample to show that a composition of maps of constant rank need not have constant rank.
\end{exercise}
